% ===== PRÉAMBULE CANONIQUE — à coller VERBATIM en tête de CHAQUE .tex =====
\documentclass[11pt,a4paper]{article}
\usepackage[utf8]{inputenc}\usepackage[T1]{fontenc}\usepackage{lmodern}
\usepackage[french]{babel}
\usepackage{amsmath,amssymb,mathtools}\usepackage{icomma}
\usepackage{siunitx}\sisetup{locale=FR}  % \qty{}{}, \si{}, \ang{}
\usepackage{xcolor}\usepackage{colortbl}
\usepackage{tikz}\usepackage{pgfplots}\pgfplotsset{compat=1.18}
\usepackage[siunitx,europeanresistors]{circuitikz} % circuits (élec.)
\usepackage[most]{tcolorbox}\usepackage{enumitem}\usepackage{array,booktabs}
\usepackage[a4paper,margin=2.2cm]{geometry}
\usetikzlibrary{arrows.meta,calc,angles,quotes,positioning,patterns,%
  backgrounds,shapes.geometric,decorations.pathmorphing,decorations.markings,intersections}
\definecolor{primary}{RGB}{222,49,99}\definecolor{primarydark}{RGB}{150,22,55}
\definecolor{defcol}{RGB}{0,114,178}\definecolor{methcol}{RGB}{52,92,148}
\definecolor{excol}{RGB}{0,135,90}\definecolor{attncol}{RGB}{200,40,55}
\tcbset{boxbase/.style={enhanced,breakable,boxrule=0.9pt,arc=2.5pt,
  left=3mm,right=3mm,top=2.3mm,bottom=2.3mm,fonttitle=\bfseries,coltitle=white}}
% --- boîtes TUTORIEL (noms EXACTS, ne pas renommer) ---
\newtcolorbox{cle}[1][]{boxbase,colback=defcol!6,colframe=defcol,
  title={\textsc{À retenir}\ifx\relax#1\relax\relax\else\textmd{ : \itshape #1}\fi}}
\newtcolorbox{methode}[1][]{boxbase,colback=methcol!7,colframe=methcol,sharp corners,
  title={\textsc{Méthode}\ifx\relax#1\relax\relax\else\textmd{ --- \itshape #1}\fi}}
\newtcolorbox{exemplebox}[1][]{boxbase,colback=excol!5,colframe=excol,
  title={\textsc{Exemple}\ifx\relax#1\relax\relax\else\textmd{ : \itshape #1}\fi}}
\newtcolorbox{piege}[1][]{boxbase,colback=attncol!6,colframe=attncol,
  title={\textsc{Piège}\ifx\relax#1\relax\relax\else\textmd{ : \itshape #1}\fi}}
\newtcolorbox{remarque}[1][]{boxbase,colback=black!4,colframe=black!45,
  title={\textsc{Remarque}\ifx\relax#1\relax\relax\else\textmd{ : \itshape #1}\fi}}
% --- boîtes EXERCICE / CORRIGÉ (noms EXACTS, auto-comptées) ---
\newtcolorbox[auto counter]{exo}[1][]{boxbase,colback=primary!4,colframe=primary,
  title={\textsc{Exercice}~\thetcbcounter\ifx\relax#1\relax\relax\else\textmd{ --- \itshape #1}\fi}}
\newtcolorbox[auto counter]{corrige}[1][]{boxbase,colback=excol!4,colframe=excol,
  title={\textsc{Corrigé}~\thetcbcounter\ifx\relax#1\relax\relax\else\textmd{ --- \itshape #1}\fi}}
\newcommand{\vv}[1]{\vec{#1}}\newcommand{\ds}{\displaystyle}
\newcommand{\R}{\mathbb{R}}\newcommand{\N}{\mathbb{N}}\newcommand{\Z}{\mathbb{Z}}
% (un \usepackage{fancyhdr}+en-tête est possible ; il est ignoré côté web)
% =========================================================================
\begin{document}
\begin{corrige}[calculer la force de Lorentz]
$\alpha=\ang{90}$ donc $\sin\alpha=1$ :
\[ F=B\,|q|\,v=0{,}5\times1{,}6\cdot10^{-19}\times10^{6}=\qty{8e-14}{\newton}. \]
\end{corrige}

\begin{corrige}[quand la force est-elle nulle ?]
\begin{enumerate}[label=\alph*)]
  \item Électron \textbf{immobile} : $v=0$, donc $F=B|q|v=\qty{0}{\newton}$. Sans mouvement, pas de
        force magnétique.
  \item Neutron : \textbf{neutre} ($q=0$), donc $F=B|q|v=\qty{0}{\newton}$, même lancé très vite.
        (Un rayon X, neutre lui aussi, n'est pas dévié non plus.)
\end{enumerate}
\end{corrige}

\begin{corrige}[sens de la force sur une charge positive]
Charge positive : on applique directement la main droite. \textbf{Pouce} vers la droite ($\vv{v}$),
\textbf{index} vers le haut ($\vv{B}$) ; le \textbf{majeur} pointe \textbf{vers toi} : la force
\textbf{sort de la feuille}, $\odot$.
\end{corrige}

\begin{corrige}[l'effet de la vitesse]
$F=B|q|v\sin\alpha$ est \textbf{proportionnelle à $v$}.
\begin{enumerate}[label=\alph*)]
  \item On double $v$ $\Rightarrow$ $F$ est \textbf{doublée} : $F=2F_0=\qty{1.6e-13}{\newton}$
        (et non quadruplée : la dépendance est \emph{linéaire}, pas quadratique).
  \item Immobile ($v=0$) $\Rightarrow$ $F=\qty{0}{\newton}$.
\end{enumerate}
\end{corrige}

\begin{corrige}[rayon de la trajectoire]
\[ R=\frac{m\,v}{|q|\,B}=\frac{1{,}67\cdot10^{-27}\times10^{6}}{1{,}6\cdot10^{-19}\times0{,}5}
    =\frac{1{,}67\cdot10^{-21}}{8\cdot10^{-20}}\approx\qty{2.1e-2}{\metre}\approx\qty{2.1}{\centi\metre}. \]
\end{corrige}

\end{document}
